
% Copyright (c) 2022 by Lars Spreng
% This work is licensed under the Creative Commons Attribution 4.0 International License. 
% To view a copy of this license, visit http://creativecommons.org/licenses/by/4.0/ or send a letter to Creative Commons, PO Box 1866, Mountain View, CA 94042, USA.

%~~~~~~~~~~~~~~~~~~~~~~~~~~~~~~~~~~~~~~~~~~~~~~~~~~~~~~~~~~~~~~~~~~~~~~~~~~~~~~
% Add your packages and commands to this file
%~~~~~~~~~~~~~~~~~~~~~~~~~~~~~~~~~~~~~~~~~~~~~~~~~~~~~~~~~~~~~~~~~~~~~~~~~~~~~~

%~~~~~~~~~~~~~~~~~~~~~~~~~~~~~~~~~~~~~~~~~~~~~~~~~~~~~~~~~~~~~~~~~~~~~~~~~~~~~~
% Fonts
% \RequirePackage{palatino} % for serif slides
% \usefonttheme{serif}
\RequirePackage[scaled]{helvet} % for sans-serif slides

\RequirePackage[utf8]{inputenc}
\RequirePackage[T1]{fontenc}

\usepackage{extsizes}
\usepackage{styles/elegantmacros}
\usepackage{adjustbox}
\usepackage{tikz}


% --- Typography -------------------------------------------------------------
% The source deck uses Gill Sans Light / SemiBold. pdfLaTeX cannot use
% arbitrary system fonts directly, so this version uses Gillius ADF, a
% Gill-Sans-inspired Type 1 font that is one of the closest TeX substitutes.
% The `scaled=0.98` adjustment gives proportions close to the reference deck.
\usepackage[scaled=0.98]{gillius}


\usefolder{styles}
\usetheme[style=lecture]{elegant}
%\usetheme[style=gold]{elegant}

\newcommand{\makepart}[1]{ % For convenience
\part{#1} \frame{\partpage}
}

%~~~~~~~~~~~~~~~~~~~~~~~~~~~~~~~~~~~~~~~~~~~~~~~~~~~~~~~~~~~~~~~~~~~~~~~~~~~~~~

%~~~~~~~~~~~~~~~~~~~~~~~~~~~~~~~~~~~~~~~~~~~~~~~~~~~~~~~~~~~~~~~~~~~~~~~~~~~~~~
% Figures
\RequirePackage{booktabs}
\RequirePackage{colortbl}
\RequirePackage{ragged2e}
\RequirePackage{schemabloc}
%\RequirePackage{natbib}
\RequirePackage{caption}
\RequirePackage{subcaption}
\RequirePackage{tabularx}
\RequirePackage{array}
\RequirePackage{multirow}
\RequirePackage[%
  natbib=true, backend=biber,%
  style=apa, isbn=false,url=false,uniquename=false%, useprefix=true%
  ]{biblatex}
\addbibresource{references.bib}
\newcolumntype{Y}{>{\centering\arraybackslash}X}

%~~~~~~~~~~~~~~~~~~~~~~~~~~~~~~~~~~~~~~~~~~~~~~~~~~~~~~~~~~~~~~~~~~~~~~~~~~~~~~

%~~~~~~~~~~~~~~~~~~~~~~~~~~~~~~~~~~~~~~~~~~~~~~~~~~~~~~~~~~~~~~~~~~~~~~~~~~~~~~
% Figures
\RequirePackage{wrapfig}
\RequirePackage{pgfplots}
\RequirePackage{graphicx}
\RequirePackage{adjustbox}
\RequirePackage{environ}
\pgfplotsset{compat=1.18}

\makeatletter
\newsavebox{\measure@tikzpicture}
\NewEnviron{scaletikzpicturetowidth}[1]{%
  \def\tikz@width{#1}%
  \def\tikzscale{1}\begin{lrbox}{\measure@tikzpicture}%
  \BODY
  \end{lrbox}%
  \pgfmathparse{#1/\wd\measure@tikzpicture}%
  \edef\tikzscale{\pgfmathresult}%
  \BODY
}
\makeatother
%~~~~~~~~~~~~~~~~~~~~~~~~~~~~~~~~~~~~~~~~~~~~~~~~~~~~~~~~~~~~~~~~~~~~~~~~~~~~~~

%~~~~~~~~~~~~~~~~~~~~~~~~~~~~~~~~~~~~~~~~~~~~~~~~~~~~~~~~~~~~~~~~~~~~~~~~~~~~~~
% Maths 
\RequirePackage{textcomp}
\RequirePackage{amsmath} 
\RequirePackage{amsthm}
\RequirePackage{mathtools}
%\RequirePackage{bbm}
%\RequirePackage{algorithm}
%\RequirePackage[osf,sc]{mathpazo}
%\RequirePackage{pifont}
%\newcommand{\xmark}{\ding{55}}%
%\numberwithin{equation}{section}
\DeclareMathOperator*{\argmax}{arg\,max}
\DeclareMathOperator*{\argmin}{arg\,min}

\setbeamertemplate{theorems}[numbered] % to number

\theoremstyle{definition}
\newtheorem{fact}{Fact}[section]
\newtheorem{examp}{Example}[section]

\theoremstyle{plain}
\newtheorem{definition}{Definition}[section]
\newtheorem{proposition}{Proposition}
\newtheorem{theorem}{Theorem}
\newtheorem{assumption}{Assumption}

\providecommand{\H}{\mathscr{H}}      
\providecommand{\E}{\mathbb{E}}
\makeatletter
\def\munderbar#1{\underline{\sbox\tw@{$#1$}\dp\tw@\z@\box\tw@}}
\makeatother

%~~~~~~~~~~~~~~~~~~~~~~~~~~~~~~~~~~~~~~~~~~~~~~~~~~~~~~~~~~~~~~~~~~~~~~~~~~~~~~


\renewcommand{\familydefault}{\sfdefault}

% --- Colors -----------------------------------------------------------------
% Dominant background sampled from the reference slides: RGB 80,103,126.
\definecolor{ubicompblue}{RGB}{80,103,126}
\definecolor{ubicompwhite}{RGB}{248,249,250}
\definecolor{ubicomplight}{RGB}{216,239,255}
\definecolor{ubicompgray}{RGB}{190,197,203}
\definecolor{ubicompdark}{RGB}{58,64,70}

% --- Global Beamer setup ----------------------------------------------------
\setbeamertemplate{navigation symbols}{}
\setbeamertemplate{headline}{}
% Page number flush right, 0.25in clear of the right paper edge and 0.1in clear
% of the bottom. The number sits in a plain \hbox to\paperwidth rather than a
% paragraph, so \parfillskip cannot compete with the \hfil and pull it back
% towards the middle. Extra box depth is what lifts the number off the bottom
% edge: with the original dp=1.2ex the digits sat 1.28pt above the paper edge,
% so dp=1.2ex-1.28pt is zero clearance and the wanted margin is added to that.
\setbeamertemplate{footline}{%
  \leavevmode%
  \hbox{%
    \begin{beamercolorbox}[wd=\paperwidth,ht=3ex,dp=\dimexpr1.2ex-1.28pt+0.1in\relax]{page number in head/foot}%
      \hbox to\paperwidth{%
        \usebeamerfont{page number in head/foot}%
        \hskip0.25in\copyright~\insertshortauthor%
        \hfil\insertframenumber\hskip0.25in%
      }%
    \end{beamercolorbox}%
  }%
}
\setbeamertemplate{items}[circle]
\setbeamersize{text margin left=0.55cm,text margin right=0.55cm}

\setbeamercolor{normal text}{fg=ubicompwhite,bg=ubicompblue}
\setbeamercolor{background canvas}{bg=ubicompblue}
\setbeamercolor{frametitle}{fg=ubicompwhite,bg=ubicompblue}
\setbeamercolor{block title}{fg=ubicompwhite,bg=ubicompblue}
\setbeamercolor{block body}{fg=ubicompwhite,bg=ubicompblue}
\setbeamercolor{structure}{fg=ubicompwhite}
\setbeamercolor{itemize item}{fg=ubicompwhite}
\setbeamercolor{itemize subitem}{fg=ubicompwhite}
\setbeamercolor{page number in head/foot}{fg=ubicompgray,bg=ubicompblue}

\setbeamerfont{normal text}{size=\fontsize{19}{23}}
\setbeamerfont{block title}{size=\fontsize{19}{23}}
\setbeamerfont{block body}{size=\fontsize{19}{23}}
\setbeamerfont{frametitle}{series=\bfseries,size=\fontsize{30}{34}}
\setbeamerfont{page number in head/foot}{size=\fontsize{8}{9}}
\setbeamerfont{itemize/enumerate body}{size=\fontsize{18}{22}}
\setbeamerfont{itemize/enumerate subbody}{size=\fontsize{16}{20}}

% --- Title page -------------------------------------------------------------
% The elegant theme puts the title and author in white panels (bg=white) and
% sets them in small caps from the structure font. Blend the panels into the
% slide background and use the same Gillius face as the frame titles.
\setbeamercolor{title}{fg=ubicompwhite,bg=ubicompblue}
\setbeamercolor{subtitle}{fg=ubicomplight,bg=ubicompblue}
\setbeamercolor{author}{fg=ubicompwhite,bg=ubicompblue}
\setbeamercolor{institute}{fg=ubicompwhite,bg=ubicompblue}
\setbeamercolor{date}{fg=ubicompgray,bg=ubicompblue}

\setbeamerfont{title}{family=\sffamily,series=\bfseries,shape=\upshape,size=\fontsize{30}{34}}
\setbeamerfont{subtitle}{family=\sffamily,series=\mdseries,shape=\upshape,size=\fontsize{20}{24}}
\setbeamerfont{author}{family=\sffamily,series=\mdseries,shape=\upshape,size=\fontsize{20}{24}}
% The theme sets these at \tiny-ish sizes; the institute carries the thanks line.
\setbeamerfont{institute}{family=\sffamily,series=\mdseries,shape=\upshape,size=\fontsize{14}{16}}
\setbeamerfont{date}{family=\sffamily,series=\mdseries,shape=\upshape,size=\fontsize{16}{18}}

% Title on the centre of the slide; author, institute (which carries the thanks
% line) and date grouped together near the foot of the slide.
% Both groups are placed absolutely with tikz overlay nodes: inside the frame's
% natural-height box \vfill collapses to zero, so glue cannot push the lower
% group down to the bottom.
\setbeamertemplate{title page}{%
  \begin{tikzpicture}[remember picture,overlay]
    \node[anchor=center,inner sep=0pt] at (current page.center) {%
      \parbox{\textwidth}{\centering\usebeamertemplate{title}}%
    };
    \node[anchor=south,inner sep=0pt] at ([yshift=0.6cm]current page.south) {%
      \parbox{\textwidth}{%
        \centering%
        \usebeamerfont{author}\usebeamercolor[fg]{author}\insertauthor\par%
        \vskip0.4em%
        \usebeamerfont{institute}\usebeamercolor[fg]{institute}\insertinstitute\par%
        \vskip0.5em%
        \usebeamerfont{date}\usebeamercolor[fg]{date}\insertdate\par%
      }%
    };
  \end{tikzpicture}%
}

% Left-aligned frame titles without decorations.
 \setbeamertemplate{frametitle}{%
   \vspace*{0.20cm}%
   \insertframetitle\par
 }

% Indent blocks to line up with other text.
% The indent is applied to the whole block (leftskip), so wrapped lines stay
% flush with the first line. Every line inside the template ends in % so no
% stray spaces leak into the block body -- those spaces used to stretch on the
% justified first line and looked like a paragraph indent.
\newlength{\blockindent}
\setlength{\blockindent}{16.4pt} % frame-title left edge - text margin - box sep

\setbeamertemplate{block begin}{%
  \par\vskip\medskipamount%
  \begin{beamercolorbox}[
    wd=\textwidth,
    leftskip=\blockindent,
    rightskip=0pt,
    sep=4pt,
    vmode
  ]{block title}%
    \usebeamerfont*{block title}\insertblocktitle%
  \end{beamercolorbox}%
  \nointerlineskip%
  \begin{beamercolorbox}[
    wd=\textwidth,
    leftskip=\blockindent,
    rightskip=0pt,
    sep=4pt,
    vmode
  ]{block body}%
    \usebeamerfont{block body}%
    \setlength{\parindent}{0pt}%
    \noindent\ignorespaces%
}

\setbeamertemplate{block end}{%
  \end{beamercolorbox}%
  \vskip\smallskipamount%
}

% \newenvironment{point}
% {%
%   \begin{tikzpicture}[remember picture,overlay]
%     \node[
%       anchor=center,
%       inner sep=0pt
%     ] at (current page.center) {%
%       \begin{adjustbox}{width=0.9\paperwidth}
%         \hbox\bgroup
% }
% {%
%         \egroup
%       \end{adjustbox}
%     };
%   \end{tikzpicture}%

% }


\NewEnviron{point}{%
  \begin{tikzpicture}[remember picture,overlay]
    \node[
      anchor=center,
      inner sep=0pt
    ] at (current page.center) {%
      \adjustbox{width=0.85\paperwidth}{%
        \strut%
        \BODY
      }%
    };
  \end{tikzpicture}%
}